\begin{abstract}
Physical changes can make a pretrained vision--language--action (VLA) model's
commands produce the wrong robot motion even when the task remains feasible.
We study embodiment adaptation from proprioceptive residuals: a calibrated
command-to-motion predictor supplies a compact physical estimate that drives
correction at the action interface while retaining the task backbone.
Healthy calibration, controlled sensitivity probes, and a specified correction
support precede deployment; subsequent updates need neither task-progress
rewards nor a healthy reference trajectory.
We compose residual calibration, partial correction support, and the actual
sampled updates with a conditional execution-recovery bound. The metric
covers the physical execution subsystem under the same supplied VLA commands;
no VLA Jacobian or full-policy contraction certificate is required. Update
holds, estimate age, clipping, realization error, and trajectory-dependent
prediction error remain explicit. Necessary contact variables must be modeled
or their effects bounded.
Across four LIBERO simulation cohorts, correction improves paired success
from 28/120 to 78/120, and from 34/120 to 68/120 under a preregistered
held-out calibration (67\% and 58\% of matched healthy performance on the
same scenarios); one robot calibration serves three VLA policies. Simulated
ALOHA and GR1 use new calibration and prior identification; from a matched
estimate an innovation-form update repairs ALOHA while adapting, so holding
is needed only where the task margin is below the update's own movement.
Four preregistered checks refuted prior readings and are reported as
primary.
A constructed servo verifies the metric and envelope; on the evaluated
robots the certificate is measured and splits, a humanoid arm certifying
while the Panda under operational-space control has no same-command
contraction, leaving its tube finite-horizon and its repair claim empirical.
\end{abstract}

\section{Introduction}
A VLA can select a useful action yet fail because the robot executes it
differently from the conditions under which the policy was trained.
Actuator offsets, reduced effectiveness, and increased friction alter the
mapping from command to motion. Those changes also alter subsequent
observations and policy decisions, creating a closed-loop physical
distribution shift. We use \emph{embodiment adaptation} to mean adaptation
of this action-to-motion interface: structured physical OOD, rather than
uncalibrated morphology transfer or missing visual and semantic competence.

Our mechanism is identification from proprioceptive prediction errors.
A healthy command-to-motion model predicts the robot response; a calibrated
sensitivity maps its residual into a small physical estimate; and that estimate
changes how the next action is realized. LIBERO's primary condition estimates six motion coordinates and corrects
three rotation coordinates -- counts of identified quantities and physical
correction directions, not of editable neural weights.

Across four LIBERO suites, correction increases paired success from
$28/120$ to $78/120$, and from $34/120$ to $68/120$ when the calibration is
identified on held-out initial states; the long-horizon suite
accounts for eight of the ten lost episodes. The same calibration serves three VLA policies. Torque
bias and heavy friction recover under a corrected fault protocol; a locked
joint is identified but not repaired. ALOHA and GR1 use new calibration and
prior fault identification, a different information budget from LIBERO's
within-episode adaptation.

\paragraph{From physical identification to recovery.}
Following selective neural fault recovery \citep{taheri2025closing}, we
connect residual calibration, adaptation, and physical execution. The nominal
physical copy follows the same supplied commands, not a rerun healthy VLA;
the execution state includes necessary robot, controller, object, and
contact variables, and a reduced servo model must bound omitted interaction
effects.

The evaluated chain is selected-coordinate estimation, applied mismatch,
scheduled execution recovery, and predictor--trajectory feedback closure.
Established same-input contraction supplies the propagation step; the
contribution is its composition with calibration error, partial support,
update holds, estimate age, and correction realization. Attenuation can
converge to a biased target, innovation can reduce truth error without
pairwise observer contraction, and the bounds preserve both mechanisms. Task
success remains empirical or requires separate margin assumptions; the
unevaluated composite extension is developed in the appendix.

\begin{figure}[t]
\centering
\resizebox{\linewidth}{!}{\begin{tikzpicture}[
  x=1cm,y=1cm,
  >=Latex,
  font=\small,
  online/.style={draw=black!70,rounded corners=2pt,align=center,
    minimum height=1.12cm,fill=blue!5},
  plant/.style={draw=black!70,rounded corners=2pt,align=center,
    minimum height=1.12cm,fill=orange!8},
  source/.style={draw=black!70,rounded corners=2pt,align=center,
    minimum height=1.12cm,fill=gray!10},
  offline/.style={draw=black!55,dashed,rounded corners=2pt,align=center,
    minimum height=.72cm,fill=gray!4,font=\scriptsize},
  signal/.style={-Latex,thick,draw=black!78},
  calibration/.style={-Latex,dashed,draw=black!55}
]
\node[source,text width=2.25cm] (vla) at (1.15,2.10)
  {Frozen VLA\\[-1pt]\textbf{unchanged}};
\node[online,text width=3.25cm] (realize) at (4.65,2.10)
  {Correction and realization\\[-1pt]
   $u_k=a_k-D\widehat f_{\tau_k}+\zeta_k$};
\node[plant,text width=2.40cm] (fault) at (8.05,2.10)
  {Modeled interface\\[-1pt]fault: $u_k^{\rm eff}=u_k+f_k$};
\node[plant,text width=2.55cm] (robot) at (11.45,2.10)
  {Low-level controller\\[-1pt]+ robot};
\node[online,text width=3.70cm] (observer) at (9.80,-.05)
  {Calibrated FIR residual\\[-1pt]
   $e_k^{\rm res}=y_k-\widehat y_k$\\[-1pt]
   $z_k=S^{-1}e_k^{\rm res}-b_{\rm cal}$};
\node[online,text width=4.10cm] (update) at (4.75,-.05)
  {Scheduled prediction-only\\[-1pt]
   update/hold: attenuation or\\[-1pt]
   innovation $\longrightarrow\widehat f_{k+1}$};
\node[offline,text width=2.50cm] (interfacecfg) at (1.15,-1.48)
  {Fixed interface:\\mask $D$;\\command limits};
\node[offline,text width=3.20cm] (updatecfg) at (5.10,-1.48)
  {Fixed update configuration:\\gain, gate, projection limits};
\node[offline,text width=3.80cm] (residualcfg) at (9.45,-1.48)
  {Fixed residual calibration:\\$H_\ell$, $S$, $b_{\rm cal}$};

\draw[signal] (vla) -- node[above] {$a_k$} (realize);
\draw[signal] (realize) -- node[above] {$u_k$} (fault);
\draw[signal] (fault) -- node[above] {$u_k^{\rm eff}$} (robot);
\draw[signal] (robot.east) -- ++(.55,0) |- node[pos=.24,right] {post-command $y_k$}
  (observer.east);
\draw[signal] (realize.south east) -- ++(0,-.35) -| node[pos=.68,above] {issued $u$-history}
  (observer.north);
\draw[signal] (observer) -- node[above] {$z_k$} (update);
\draw[signal] (update.north) -- ++(0,.36) -| node[pos=.73,left] {available snapshot
  $\widehat f_{\tau_k}$} (realize.south west);
\draw[calibration] (interfacecfg.north east) -- (realize.west);
\draw[calibration] (updatecfg.north) -- (update.south);
\draw[calibration] (residualcfg.north) -- (observer.south);

\end{tikzpicture}
}
\caption{Evaluated prediction-only execution adapter. The frozen VLA emits
$a_k$; the available estimate snapshot is masked and realized before the
modeled additive interface fault $f_k$. After execution, the calibrated FIR
residual $z_k$ drives the scheduled attenuation or innovation update. The
task--observation loop is omitted; the theorem conditions on the supplied
$a_k$. No nominal rollout, geodesic signal, task reward, or VLA update is
used online (Appendix~\ref{app:figure-record}).}
\label{fig:framework}
\end{figure}
